\documentclass{article}
\usepackage{graphicx} % Required for inserting images
\usepackage[a4paper, margin=1in]{geometry}

\title{Project proposal template}
\author{Your names}
\date{Date of submission}

\begin{document}

\maketitle

\section*{Guidelines for Project Proposal LMP2392H (2 page max)}

\section{Title}
Specific and Descriptive. E.g. Predicting Post-Operative Complications Using EHR Time-Series Data

\section{Team}
Students and the clinical/senior lead. You may want to include your expertise as it applies to the project

\section{Problem Statement \& Clinical Motivation}
Clearly describe the clinical or health-system problem being addressed, who it affects, and why it
matters in practice. Ground the problem in a real clinical workflow and avoid vague objectives.

\section{Relevant Work}
Include at least some of the relevant existing AI work that solves this or a similar problem. 

\section{Data \& Setting}
Describe the data type(s), source (public, synthetic, or hypothetical), approximate sample size, and key
variables. Evaluation focuses on plausibility, not actual data access.

\section{Methods}
Explain the modeling plan, rationale for the approach, and baselines. Please include an alternative plan if your initial idea doesn’t work. Novelty is optional; justification and appropriateness are required. 

\section{Evaluation Strategy}
Define primary metrics and validation strategy. Justify metric choice in clinical context and consider
subgroup or fairness analysis where relevant.

\section{Deployment Context \& Limitations}
Describe where the system fits in the clinical workflow, its role (decision-support vs automation), and
key limitations or risks such as data shift, bias, interpretability, or adoption barriers.

\section{Project Scope \& Deliverables}
Clearly state expected deliverables and explicitly list what is out of scope.


\paragraph{\textbf{Notes}:}
\begin{itemize}
    \item You may want to include citations to clinically ground your project in available AI solutions (esp for the related work section)
    \item Make sure that the project is doable within a semester (feasibility over ambition, but creativity is always welcome and encouraged!)
    \item Proposals will be evaluated on clinical grounding, clarity, methodological appropriateness, evaluation soundness, and awareness of deployment and limitations.
\end{itemize}

\end{document}
